\documentclass[10pt,a4paper]{amsart}
\usepackage[margin=19mm]{geometry}
\usepackage{amsmath,amssymb,mathtools}
\usepackage[scr=rsfs,cal=euler]{mathalfa}
\usepackage{xfrac}
\usepackage[unicode,hidelinks,bookmarks=false]{hyperref}
\newcommand{\ie}{i.e.\ }
\newcommand{\cf}{cf.\ }
\newcommand{\resp}{respectively\ }
\newcommand{\Subsection}{\subsection}
\providecommand{\nobreakdash}{\nobreak\hbox{-}}
\setlength{\emergencystretch}{2em}
\multlinegap0pt
\usepackage{amssymb}
\newtheorem{lemma}{Lemma}
\title{Toda flow with unbounded initial data}
\author{Shinichi Kotani, Jiahao Xu, Shuo Zhang}
\date{Source: arXiv:2604.05434v1 (CC BY 4.0)}
\begin{document}
\maketitle

\noindent\emph{Source and licence.} Toda flow with unbounded initial data. Version arXiv:2604.05434v1 (CC BY 4.0), distributed by arXiv under the Creative Commons Attribution 4.0 International licence, http://creativecommons.org/licenses/by/4.0/, as recorded in the arXiv licence metadata for this exact version and shown on the abstract page https://arxiv.org/abs/2604.05434v1. Full licence notice: \texttt{licenses/arxiv-2604.05434-CC-BY-4.0.txt}.\par\smallskip

\noindent\textit{Editorial scope.} Counted unit: the article's lemma l2-8 (source0.tex lines 1792--1800). The article's complete original proof of it is delivered with it (source0.tex lines 1802--1950, one \texttt{proof} environment of 149 lines). No mathematics is written, repaired or re-derived: the statement and the proof are the article's own lines, in the article's own notation, and no retained line is substituted or re-flowed.  Retained uncounted as the source-local context the proof needs: lines 1345--1363.  Nothing was omitted from any retained span, so the package carries no omission record.  No retained line contains a citation, so the wrapper carries no bibliography and no bibliography entry.  No retained line contains a figure environment, an image-inclusion command or drawing code, so no image is delivered and the package declares no asset dependency.  Editorial formatting only, in addition: an AMS \texttt{amsart} preamble that restates the article's own declarations and macros used by the retained lines, namely the environments \texttt{lemma} on separate counters, together with the standard packages those lines need.  The retained environments print in this document as Lemma 1 in order of appearance, whereas the article prints the same items with the numbering of its own full document.  The complete native source of the article as distributed in the arXiv e-print for this exact version is bundled unchanged as \texttt{sources/197956/source0.tex}.
\begin{lemma}
\label{l2-4}For $g_{1}$, $g_{2}\in\Gamma_{N}\left(  D_{+}\right)  $ and any
$c>c^{\prime}\equiv\max_{j}\left\{  \mathrm{\delta}_{D_{+}}\left(
g_{j}\right)  \right\}  $,
\[
\mathrm{d}_{c}\left(  g_{1},g_{2}\right)  \leq\sqrt{\int_{C^{2}}\left\vert
\dfrac{G_{c^{\prime}}\left(  z,\lambda\right)  }{\lambda-z}\right\vert
^{2}\left\vert \rho_{c}\left(  \lambda\right)  \right\vert ^{-2}\left\vert
\rho_{c}\left(  z\right)  \right\vert ^{-2}\left\vert d\lambda\right\vert
\left\vert dz\right\vert }<\infty
\]
is valid. Especially, $\mathrm{d}_{c}\left(  g,1\right)  <\infty$ holds for
$g\in\Gamma_{N}\left(  D_{+}\right)  $, if $c>\mathrm{\delta}_{D_{+}}\left(
g\right)  $. Moreover, for a sequence $g_{n}\in\Gamma_{N}\left(  D_{+}\right)
$, suppose $\sup_{n\geq1}\mathrm{\delta}_{D_{+}}\left(  g_{n}\right)  <c$ and
$g_{n}\left(  \lambda\right)  \rightarrow g\left(  \lambda\right)  $ for a.e.
$\lambda\in C$. Then, $\mathrm{d}_{c}\left(  g_{n},g\right)  \rightarrow0$ as
$n\rightarrow\infty$ holds.\newline
\end{lemma}

\begin{lemma}
\label{l2-8}For $g\in\Gamma_{N}\left(  D_{+}^{j}\right)  $ there exists
rational functions $\left\{  r_{n}\right\}  _{n\geq1}\subset\Gamma_{N}\left(
D_{+}^{j}\right)  $ such that for sufficiently large $c$%
\[
\lim_{n\rightarrow\infty}\mathrm{d}_{c}\left(  r_{n},g\right)  =0\text{.}%
\]

\end{lemma}

\begin{proof}
Let $g\in\Gamma_{N}\left(  D_{+}^{1}\right)  $ and assume $g(0)=1$ without
loss of generality. The proof consists of 3 steps.

\textbf{Step} \textbf{1}. Set%
\[
h\left(  z\right)  =%
%TCIMACRO{\dint _{0}^{z}}%
%BeginExpansion
{\displaystyle\int_{0}^{z}}
%EndExpansion
\dfrac{g^{\prime}\left(  \zeta\right)  }{g\left(  \zeta\right)  }d\zeta\text{,
\ }g_{n}\left(  z\right)  =\exp\left(  \int_{0}^{z}\dfrac{h^{\prime}\left(
\lambda\right)  }{1+\lambda^{2}/n}d\lambda\right)  \in\Gamma_{N-1}\left(
D_{+}^{1}\right)  \text{.}%
\]
Then, $g_{n}\left(  z\right)  \rightarrow$ $g\left(  z\right)  $ on $\left(
D_{+}^{\prime}\right)  ^{1}$ and%
\[
\left\vert \dfrac{g_{n}^{\prime}\left(  z\right)  }{g_{n}\left(  z\right)
}\right\vert =\dfrac{\left\vert h^{\prime}\left(  z\right)  \right\vert
}{\left\vert 1+z^{2}/n\right\vert }\leq c_{1}\left\vert z\right\vert
^{N-1}\text{ \ on }\left(  D_{+}^{\prime}\right)  ^{1}\text{ if }\left\vert
z\right\vert \geq1\text{,}%
\]
which implies $\sup_{n\geq1}\mathrm{\delta}_{D_{+}^{1}}\left(  g_{n}\right)
<\infty$, and $\lim_{n\rightarrow\infty}\mathrm{d}_{c}\left(  g_{n},g\right)
=0$ for sufficiently large $c$ due to Lemma \ref{l2-4}.

\textbf{Step} \textbf{2}. Assume $g\in\Gamma_{N-1}\left(  D_{+}^{1}\right)  $
and set $h\left(  z\right)  =\int_{0}^{z}g^{\prime}\left(  \zeta\right)
/g\left(  \zeta\right)  d\zeta$ again. We approximate $h(z)$ by rational
functions. Let $f$ be an integrable function on the curve $C_{1}^{\prime}$.
Then, for $M>0$, $u\left(  z\right)  \equiv\dfrac{1}{2\pi i}\int_{C^{\prime}%
}\dfrac{1}{\lambda-z}f\left(  \lambda\right)  d\lambda$ has a decomposition:%
\begin{equation}
u\left(  z\right)  =\dfrac{1}{2\pi i}\int_{C_{1}^{\prime}:\left\vert
\lambda\right\vert \leq M}\dfrac{1}{\lambda-z}f\left(  \lambda\right)
d\lambda+\dfrac{1}{2\pi i}\int_{C_{1}^{\prime}:\left\vert \lambda\right\vert
>M}\dfrac{1}{\lambda-z}f\left(  \lambda\right)  d\lambda\text{.} \label{2-16}%
\end{equation}
Noting $\mathrm{dist}\left(  C_{1},C_{1}^{\prime}\right)  >0$, one has%
\[
\sup_{z\in D_{+}^{1}}\left\vert \dfrac{1}{2\pi i}\int_{C_{1}^{\prime
}:\left\vert \lambda\right\vert >M}\dfrac{1}{\lambda-z}f\left(  \lambda
\right)  d\lambda\right\vert \leq\dfrac{1}{2\pi\mathrm{dist}\left(
C_{1},C_{1}^{\prime}\right)  }\int_{C_{1}^{\prime}:\left\vert \lambda
\right\vert >M}\left\vert f\left(  \lambda\right)  \right\vert \left\vert
d\lambda\right\vert \text{,}%
\]
which can be made small as much as possible by choosing a large $M$. On the
other hand, the first term of (\ref{2-16}) is an integral on a compact set. If
a compact set $K$ of $\mathbb{C}$ has diameter $\varepsilon$, choosing $b\in
K$, the expansion:%
\[
\dfrac{1}{2\pi i}\int_{K}\dfrac{1}{\lambda-z}f\left(  \lambda\right)
d\lambda=\left(  b-z\right)  ^{-1}\sum_{j\geq0}\left(  b-z\right)  ^{-j}%
\dfrac{1}{2\pi i}\int_{K}\left(  b-\lambda\right)  ^{j}f\left(  \lambda
\right)  d\lambda
\]
converges uniformly on the region $\left\vert b-z\right\vert \geq2\varepsilon
$. A Cauchy integral on a general compact set $K$ can be approximated by
rational functions uniformly on the $2\varepsilon$-neighborhood of $K$ by
covering $K$ by a finite number of $\varepsilon$-disks. Hence, $u(z)$ can be
approximated by rational functions with no poles in $D_{+}^{1}$ uniformly on
$D_{+}^{1}$. Clearly, the derivative $u^{\prime}(z)$ also can be approximated
by the derivatives of the rational functions uniformly. Applying this argument
to $f(\lambda)=\lambda^{-1}h\left(  \lambda\right)  \left(  \lambda-b\right)
^{-N}\in L^{1}\left(  C_{1}^{\prime}\right)  $ with $b\in\mathbb{C}%
\diagdown\overline{D_{+}^{1}}$, one sees that $f\left(  z\right)  $ can be
approximated uniformly by rational functions $s_{k}(z)$ with no poles on
$\overline{D_{+}^{1}}$. Note $s_{k}(z)=O\left(  z^{-1}\right)  $ as
$z\rightarrow\infty$. In this case, since $h^{\prime}\left(  z\right)
=O\left(  z^{N-1}\right)  $, one has $s_{k}^{\prime}(z)=O\left(
z^{-2}\right)  $ as $z\rightarrow\infty$ uniformly in $k\geq1$ as well.

\textbf{Step 3}. Set
\[
g_{n}\left(  z\right)  =\left(  1+\dfrac{z^{2N}-h\left(  z\right)  }%
{n}\right)  ^{-n}\left(  1+\dfrac{z^{2N}}{n}\right)  ^{n}\text{,}%
\]
and define rational functions:%
\[
r_{n}\left(  z\right)  =\left(  1+\dfrac{z^{2N}-h_{n}\left(  z\right)  }%
{n}\right)  ^{-n}\left(  1+\dfrac{z^{2N}}{n}\right)  ^{n}\text{ with }%
h_{n}\left(  z\right)  =z\left(  z-b\right)  ^{N}s_{n}\left(  z\right)
\text{.}%
\]
It is clear that $\lim_{n\rightarrow\infty}r_{n}\left(  z\right)  =e^{h\left(
z\right)  }=g(z)$ uniformly on each compact set of $\overline{D}_{+}^{1}$. We
show $\sup_{n\geq n_{1}}\mathrm{\delta}_{D_{+}^{1}}\left(  r_{n}\right)
<\infty$ for some $n_{1}\geq1$. Observe%
\[
\dfrac{r_{n}^{\prime}\left(  z\right)  }{r_{n}\left(  z\right)  }=\dfrac
{h_{n}^{\prime}(z)}{1+\left(  z^{2N}-h_{n}\left(  z\right)  \right)
/n}-\dfrac{2Nz^{2N}/n}{1+z^{2N}/n}\dfrac{z^{-1}h_{n}\left(  z\right)
}{1+\left(  z^{2N}-h_{n}\left(  z\right)  \right)  /n}\text{.}%
\]
For $M>0$ set%
\[
\left\{
\begin{array}
[c]{l}%
c_{1}=\inf_{z\in D_{+}^{1},\left\vert \operatorname{Re}z\right\vert
>M}\operatorname{Re}\left(  z^{2N}-h_{n}\left(  z\right)  \right)  \left(
\operatorname{Re}z\right)  ^{-2N}\\
c_{2}=\sup_{z\in D_{+}^{1},\left\vert \operatorname{Re}z\right\vert \leq
M}\left\vert z^{2N}-h_{n}\left(  z\right)  \right\vert
\end{array}
\right.  \text{.}%
\]
We have $c_{1}>0$ for sufficiently large $M$, since $h_{n}\left(  z\right)
=O\left(  z^{N}\right)  $ as $z\rightarrow\infty$ uniformly on $n$ and
$\sup_{z\in D_{+}^{1}}\left\vert \operatorname{Im}z\right\vert <\infty$. Thus,
we have%
\[
\left\vert 1+\dfrac{z^{2N}-h_{n}\left(  z\right)  }{n}\right\vert \geq\left\{
%
\begin{array}
[c]{ll}%
1+c_{1}\left(  \operatorname{Re}z\right)  ^{2N}/n & \text{if \ }\left\vert
\operatorname{Re}z\right\vert >M\text{, }z\in D_{+}^{1}\\
1-c_{2}/n & \text{if \ }\left\vert \operatorname{Re}z\right\vert \leq M\text{,
}z\in D_{+}^{1}%
\end{array}
\right.  \text{.}%
\]
Similarly, one has%
\[
\left\vert \dfrac{z^{2N}/n}{1+z^{2N}/n}\right\vert \leq c_{3}\dfrac{\left(
\operatorname{Re}z\right)  ^{2N}/n}{1+\left(  \operatorname{Re}z\right)
^{2N}/n}\leq c_{3}%
\]
with some constant $c_{3}$. Consequently,%
\[
\left\vert \dfrac{r_{n}^{\prime}\left(  z\right)  }{r_{n}\left(  z\right)
}\right\vert \leq c_{4}\left(  \left\vert h_{n}^{\prime}(z)\right\vert
+\left\vert z\right\vert ^{-1}\left\vert h_{n}\left(  z\right)  \right\vert
\right)
\]
for another constant $c_{4}$, which implies $\sup_{n\geq n_{1}}\mathrm{\delta
}_{D_{+}^{1}}\left(  r_{n}\right)  <\infty$ (note $h_{n}(0)=0$ and
$h_{n}\left(  z\right)  =O\left(  z^{N}\right)  $, $h_{n}^{\prime}\left(
z\right)  =O\left(  z^{N-1}\right)  $ uniformly in $n$) for $n_{1}>c_{2}$, and
Lemma \ref{l2-4} yields $\mathrm{d}_{c}\left(  r_{n},g\right)  \rightarrow0$
for a sufficiently large $c>0$. The case $g\in\Gamma_{N}\left(  D_{+}%
^{2}\right)  $ can be treated by applying the above argument to $\widetilde
{g}\in\Gamma_{N}\left(  D_{+}^{1}\right)  $.\bigskip
\end{proof}

\end{document}
